\documentclass[12pt]{article}
\usepackage[utf8]{inputenc}
\usepackage[T1]{fontenc}
\usepackage{hyperref}
\title{Tetun Dili Language Resources for K-12 Curriculum Development}
\author{Compiled by Claude Sonnet 5 (claude-sonnet-5) for Andr\'as Kornai}
\date{}
\begin{document}
\maketitle
\section{Introduction}
Tetun Dili (ISO 639-3 tdt) is, with Portuguese, an official language of Timor-Leste, centred on Dili, and a language of instruction. This search found reference and learner materials but no Ministry of Education school textbook series; the school materials remain to be identified.
\section{Primary Educational Resources}
\begin{itemize}
\item \emph{Mai Aprende Tetum ho Ha'u} (Adelaide Lopes), a learner text listed in the ANU Tetum course: \url{https://programsandcourses.anu.edu.au/course/TETM1002/First%20Semester/2245}
\end{itemize}
These are for adult learners, not K-12 pupils.
\section{Linguistic and Reference Materials}
\begin{itemize}
\item \emph{Tetun Dili: A Grammar of an East Timorese Language} (Williams-van Klinken, Hajek, Nordlinger 2002), review: \url{https://linguistlist.org/issues/12/3024/}
\item \emph{Tetum Reference Grammar} (Hull and Eccles) and Hull's \emph{Short English-Tetum Dictionary}, in the ETAN catalogue: \url{https://www.etan.org/resource/learntetum.htm}
\item \emph{Tetun Word-Finder} (Williams-van Klinken, 2008): \url{https://mailman.anu.edu.au/pipermail/easttimorstudies/2009-March/000614.html}
\item Online Tetun Dili dictionary on Webonary: \url{https://www.webonary.org/tetun-dili/overview/introduction/}
\end{itemize}
\section{Religious and Community Resources}
None found in this search.
\section{Digital and Audio Resources}
The Webonary dictionary is the only digital resource found.
\section{Curriculum Development Context}
Tetun is already used in education, but this search did not locate official curricula or textbooks. Further search of Timor-Leste Ministry of Education materials and mother-tongue-based programmes is recommended.
\end{document}